\section{Database Fields and Corpus Summary}
\label{app:database}

The evidence database records bibliographic identifiers and version relations; inclusion role; primary and secondary clinical applications; capability level; modality; disease and organ system; clinical task; state representation; action or intervention; transition model; outcome or reward; counterfactual and planning support; datasets and access conditions; code and model availability; validation setting; source locations; confidence; and adjudication status. Architecture remains a provisional field in this draft. Quantitative architecture analysis will be added only after a stratified coding study establishes stable definitions and resolves disagreements across model families.

The corpus summary tables are descriptive rather than measures of clinical maturity. Capability counts report implemented task form. Application counts describe the primary clinical setting. Resource counts describe what was publicly verifiable at the time of review. None of these dimensions substitutes for causal validity, external transportability, prospective evaluation, or regulatory readiness.

\ifworkingdraft
\paragraph{Living resource.} A maintained reading list is available at \url{https://github.com/lanqz7766/awesome-medical-world-models}.
\fi

\begin{table*}[p]
\caption{Representative works across six clinical applications. Capability, application, disease family, modality, task, and public resource status are independent fields; digital twin is treated as a system-form overlay.}
\label{tab:track-works}
\centering
\scriptsize
\setlength{\tabcolsep}{3pt}
\renewcommand{\arraystretch}{1.02}
\begin{tabularx}{\textwidth}{@{}p{3.45cm}ccp{2.45cm}p{1.9cm}Xp{1.1cm}@{}}
\toprule
\textbf{Representative work} & \textbf{Year} & \textbf{Level} & \textbf{Disease family} & \textbf{Modality} & \textbf{Evaluated task} & \textbf{Code} \\
\midrule
\multicolumn{7}{@{}l}{\textit{Medical imaging}} \\
CheXWorld: Exploring Image World Modeling for Radiograph Representation ... & 2025 & 1 & general or multidisease & radiograph & representation & \cgood yes \\
Xray2Xray: World Model from Chest X-rays with Volumetric ... & 2025 & 3 & general or multidisease & radiograph & action conditioned simulation & \cgap no* \\
MRI Contrast Enhancement Kinetics World Model & 2026 & 2 & oncology & MRI & future observation synthesis & \cgood yes \\
\addlinespace
\multicolumn{7}{@{}l}{\textit{Longitudinal progression}} \\
Forecasting individualized progression of Alzheimer's disease using structural ... & 2026 & 2 & neurologic disease & MRI & future observation synthesis & \cgap no* \\
Treatment-aware Diffusion Probabilistic Model for Longitudinal MRI Generation ... & 2025 & 3 & oncology & MRI & treatment conditioned simulation & \cgap no* \\
Explicit Temporal Embedding in Deep Generative Latent Models ... & 2023 & 2 & oncology & CT & future observation synthesis & \cgood yes \\
\addlinespace
\multicolumn{7}{@{}l}{\textit{EHR trajectories}} \\
EHRWorld: A Patient-Centric Medical World Model for Long-Horizon ... & 2026 & 3 & general or multidisease & EHR & future forecasting & \cgap no* \\
ChronoMedicalWorld: A Medical World Model for Learning Patient ... & 2026 & 3 & renal urologic disease & EHR & disease progression & \cgap no* \\
The Patient is not a Moving Document: A ... & 2026 & 2 & oncology & EHR & future forecasting & \cgap no* \\
\addlinespace
\multicolumn{7}{@{}l}{\textit{Digital twins}} \\
A Framework for the generation of digital twins ... & 2021 & 1 & not disease specific & MRI & representation & \cgap no* \\
Development of Digital Twins to Optimize Trauma Surgery ... & 2021 & 5 & musculoskeletal disease & CT & planning decision support & \cgap no* \\
A Novel Cloud-Based Framework for the Elderly Healthcare ... & 2019 & 3 & general or multidisease & wearable sensor & operations resource simulation & \cgap no* \\
\addlinespace
\multicolumn{7}{@{}l}{\textit{Treatment and causal modeling}} \\
ECG-WM: A Physiology-Informed ECG World Model for Clinical ... & 2026 & 5 & cardiovascular disease & physiology & planning decision support & \cgap no* \\
Learning to Treat Hypotensive Episodes in Sepsis Patients ... & 2021 & 4 & critical illness & EHR & treatment policy learning & \cgap no* \\
auton-survival: an Open-Source Package for Regression, Counterfactual Estimation, ... & 2022 & 4 & oncology & demographics clinical attributes & treatment response & \cgood yes \\
\addlinespace
\multicolumn{7}{@{}l}{\textit{Surgery, robotics, and physiology}} \\
SWoMo: Neuro-Symbolic World Model for Cataract Surgery Simulation & 2026 & 3 & ophthalmic disease & surgical video/robot & simulation training & \cgood yes \\
EchoWorld: Learning Motion-Aware World Models for Echocardiography Probe ... & 2025 & 3 & not disease specific & US/echo & acquisition guidance & \cgood yes \\
SAW: Toward a Surgical Action World Model via ... & 2026 & 3 & not disease specific & surgical video/robot & simulation training & \cgap no* \\
\addlinespace
\bottomrule
\end{tabularx}
\end{table*}

\begin{table}[t]
\caption{Public-resource status at the review date. Code categories describe publicly verifiable artifacts; ``not found'' does not imply that private code is absent. Unresolved full texts are excluded from the primary scientific distributions.}
\label{tab:openness}
\centering
\small
\begin{tabular}{p{1.7cm}p{7.0cm}r}
\toprule
\textbf{Axis} & \textbf{Category} & \textbf{Count} \\
\midrule
Code & paper-linked public code verified & 111 \\
Code & public but partial/related or without weights & 6 \\
Code & claimed, announced, by request, or pending & 19 \\
Code & not found, unavailable, proprietary, or unreleased & 96 \\
Code & not applicable to review/perspective & 68 \\
Full text & sufficient evidence for paper-level review & 300 \\
Full text & unresolved and excluded from primary distributions & 23 \\
\bottomrule
\end{tabular}
\end{table}

